\section{Finite-step simulation and numerical evidence}\label{sec:numerics}

The drift restriction of Section~\ref{sec:drift} admits two numerically distinct
implementations.  The first integrates the log-DLV SDE directly.  It is the
closest implementation of the continuous-time model, but a finite time step is
only approximately martingale preserving.  The second advances the probability
density itself through a finite set of symmetric branches.  It is an exact
one-step martingale construction on the modeled finite support, although it is
not a claim about a continuous-strike or unbounded-support limit.  We report
both because their numerical guarantees differ.

All experiments in this section use four weekly expiries and fixed cash strikes
within each weekly transition.  ``DLVsq'' denotes the positive discrete local
variance $\dlv$.  The final density-branch audit uses the discrete factor $X$
alone.  Earlier development and long-horizon audits also multiply $X$ by an
independent mean-one Black factor $Y$ with annual volatility 10\%; the
conditioning identities below apply to both, using the Black semigroup for
$Y$.

\subsection{Direct integration of the DLV diffusion}\label{sec:numerics-sde}

For the direct SDE experiment we work in log excess-DLV coordinates
\[
  z^{j,i}=\log(\dlv^{j,i}-\dlv_{\min}).
\]
The volatility of $z$ is the rank-25 residual covariance estimated from 703
weekly historical transitions through 31 December 2024.  At every substep the
density volatility is obtained from \eqref{eq:volsolve}, and the drift of
$\dlv$ is then computed from \eqref{eq:pointwise}, including the It\^o term in
log coordinates.  Each week is integrated by ten stochastic-Heun substeps,
using the same Gaussian increment in predictor and corrector.  At the weekly
endpoint spot is drawn from the front density and the expired slab is removed.

The audit starts from historical state 44 (8 March 2022), uses 200 interior
strikes, 100,000 paths arranged as 50,000 antithetic pairs, and follows a
shrinking tower for four weeks without re-gridding.  Fifty antithetic pairs are
aggregated into each of 1,000 independent randomized-stratification blocks.
There are 800, 600, 400, and 200 live fixed-contract option returns at the four
trade dates, hence 2,000 coordinates in the joint test.

For this experiment the blockwise coordinate statistics are supplemented by a
10,000-draw shared-sign multiplier bootstrap.  If $t_k$ is the usual block
t-statistic, the two simultaneous statistics are
\[
  T_{\max}=\max_k|t_k|,
  \qquad
  T_{\mathrm{RMS}}=\left(d^{-1}\sum_{k=1}^{d}t_k^2\right)^{1/2}.
\]
The reported decision value is
$p_{\mathrm{comb}}=\min\{1,2\min(p_{\max},p_{\mathrm{RMS}})\}$.
The shared multiplier for all coordinates of a block preserves their observed
cross-option dependence.

\begin{table}[htbp]
\centering
\small
\begin{tabular}{lrrrrl}
\toprule
Audit & Active & $\max|t|$ & RMS $t$ & $p_{\mathrm{comb}}$ & 5\% decision\\
\midrule
One-step option returns       & 2,000 & 3.334 & 1.109 & 0.2746 & do not reject\\
Terminal option returns       & 2,000 & 3.334 & 1.044 & 0.3210 & do not reject\\
Terminal option price changes & 2,000 & 2.934 & 1.156 & 0.3500 & do not reject\\
Integrated-DLV control        & 1,200 & 1.950 & 0.878 & 1.0000 & do not reject\\
Weekly spot returns           & 4     & 2.431 & 1.561 & 0.0876 & do not reject\\
\bottomrule
\end{tabular}
\caption{Martingale diagnostics for the ten-substep Heun implementation.  The
800 omitted integrated-DLV coordinates have effectively zero block variance.
The primary one-step audit has separate multiplier-bootstrap values
$p_{\max}=0.1373$ and $p_{\mathrm{RMS}}=0.2465$; its simultaneous 5\% max-$t$
cutoff is 4.235.}
\label{tab:sde-martingale}
\end{table}

The largest density-mass residual was $9.33\times10^{-14}$ and the largest
conditional-forward residual was $8.93\times10^{-14}$.  The minimum density was
zero, due to finite-support point masses rather than a negative density.  The
simulation took 72.1 seconds, or approximately 1,388 paths per second.  These
results support the discretization, but they do not turn an infinitesimal drift
restriction into an exact one-week transition.

\subsection{An exact finite-step density construction}\label{sec:density-step}

We next construct the one-week law in density coordinates.  Let
$\varepsilon_b\in\{-1,+1\}^{d}$ be the $b$-th Rademacher branch and let
$\pi_b=2^{-d}$, with $d=4$ in the experiment.  Write $\Gamma^j$ for the
state-dependent loading of density $q^j$ implied by the chosen DLVsq loading.
The branch marginal is
\begin{equation}\label{eq:density-branch}
  q^{j,b}=q^j+\sqrt h\,\Gamma^j\varepsilon_b,
  \qquad b=1,\ldots,2^d .
\end{equation}
Since
$\sum_b\pi_b\varepsilon_b=0$ and
$\sum_b\pi_b\varepsilon_b\varepsilon_b^{\!\top}=I_d$,
\begin{equation}\label{eq:density-branch-mean}
  \sum_b\pi_bq^{j,b}=q^j
\end{equation}
holds exactly, before Monte Carlo sampling.  DLVsq is then recovered
algebraically from each $q^{j,b}$ by Proposition~\ref{prop:codebook}; it is not
advanced by an Euler step or fitted by a nonlinear optimizer.  A branch is
accepted only if its masses are non-negative and the recovered DLVsq lies in
its configured open bounds.  The implementation checks mass, mean, decoder
consistency, and DLV bounds at every step.

In the implemented weekly law the front density is branch invariant,
$q^{1,b}=q^1$.  Independently of $b$, a front node $I=i$ is drawn with
probability $q^1_i$.  For every surviving expiry $j\ge2$, let
$R^{b,j}_{ir}$ be the conditional probability of node $r$ given branch $b$ and
front node $i$.  The SANOS transition supplies rows satisfying
\begin{align}
  \sum_rR^{b,j}_{ir}&=1,
  &\sum_rx_rR^{b,j}_{ir}&=x_i,
  &\sum_iq^1_iR^{b,j}_{ir}&=q^{j,b}_r .
  \label{eq:conditional-density-identities}
\end{align}
The first two identities make every row a martingale transition; the third
recovers the branch marginal.  Together with
\eqref{eq:density-branch-mean}, they imply the tower property for every payoff
represented on the finite support.

This construction is exact on the finite model grid.  It is not asserted to be
the unique risk-neutral law, and the particular independent sampling of $b$
and $I$ does not generate contemporaneous correlation between the raw surface
branch and front-node innovation.  Realized density-weighted factor covariance
nevertheless matched the intended empirical covariance closely: in the
independent all-state confirmation the largest relative error was 1.87\% and
the RMS error was 0.55\%.

\subsection{Conditioning and re-gridding}\label{sec:numerics-regrid}

Suppose current cash spot is $s$ and the revealed front node is $x_i$, so the
next $X$ spot is $s x_i$.  Define the next relative support by
\begin{equation}\label{eq:exact-regrid}
  x_r^{\prime(i)}=\frac{x_r}{x_i}.
\end{equation}
Then $(s x_i)x_r^{\prime(i)}=s x_r$ path by path.  Consequently the price of a
surviving fixed-cash call is identical on the inherited and re-gridded
coordinates,
\[
  C_{t+h}^{b,i,j}(K)
  =\sum_rR^{b,j}_{ir}(s x_r-K)^+.
\]
Averaging first over $i$ and then over $b$ gives
\begin{align*}
  \E_t[C_{t+h}^{j}(K)]
  &=\sum_b\pi_b\sum_iq^1_i\sum_rR^{b,j}_{ir}(s x_r-K)^+\\
  &=\sum_rq^j_r(sx_r-K)^+=C_t^j(K).
\end{align*}
Thus re-gridding is an exact coordinate change for all surviving fixed-cash
claims, not an interpolation.  With an independent Black factor $Y$, the same
proof conditions on $Y$ and then applies its mean-one Black semigroup.  A newly
issued far expiry is different: there is no pre-existing fifth-week claim whose
price it must preserve.

At each weekly roll, the old 2W--4W densities become the new floating 1W--3W
densities through \eqref{eq:conditional-density-identities} and
\eqref{eq:exact-regrid}.  Only the new 4W DLVsq curve is supplied statistically.
The retained boundary model is a 25-factor PCA followed by an unrestricted
weekly VAR(1) in factor scores.  It uses 703 training pairs through 2024,
explains 99.9912\% of the 4W curve variation, and has spectral radius 0.9089.
On 155 untouched 2025 pairs its raw-coordinate RMSE is 0.5902, against 0.6526
for persistence, a 9.56\% improvement.  SANOS constructs the new 4W density as
a static martingale extension of the inherited 3W density.

\subsection{Martingale audit for the exact density step}\label{sec:density-tests}

Adjacent simulated paths form an antithetic pair.  For pair $a=1,\ldots,n$ and
audited return coordinate $k$, let $R_{ak}$ be the pair mean, $\bar R$ the
vector of sample means, and $\widehat\Sigma$ their sample covariance.  We
report
\begin{equation}\label{eq:density-wald}
  z_k=\frac{\sqrt n\,\bar R_k}{\widehat\sigma_k},
  \qquad
  J=n\bar R^{\!\top}\widehat\Sigma^{+}\bar R .
\end{equation}
The pseudo-inverse is formed in correlation space after discarding eigenmodes
at or below 0.01.  If $\nu$ modes remain, the reference value is
$J\sim\chi^2_\nu$ under the joint zero-mean null.  A small upper-tail $p$-value
rejects zero drift; a large value means only that the observed drift is
compatible with Monte Carlo noise.  We use 0.05 for one look, 0.0125 for four
weekly looks, and 0.003125 for 16 looks.  Because the options are highly
collinear and the retained modes are estimated from the sample, the exact
density-tower identities and matched cash/re-grid prices are stronger
structural diagnostics than an isolated asymptotic $p$-value.

The canonical X-only audit used all 891 certified historical states, 100,000
paths (50,000 independent antithetic pairs), 11 ITM call/put strikes per
expiry, four one-week steps, 64-bit arithmetic, and 10\% annualized log-DLVsq
volatility.  Each state received 56 or 57 independent pairs.  The audited book
contains spot and the surviving original fixed-cash claims.

\begin{table}[htbp]
\centering
\small
\begin{tabular}{rrrrrr}
\toprule
Week & Worst $|z|$ & Max option mean (bp) & $J$ & $\nu$ & $p$\\
\midrule
1 & 1.796 & 0.497 & 10.980 & 7 & 0.139\\
2 & 2.143 & 0.863 &  7.228 & 7 & 0.406\\
3 & 1.371 & 0.696 &  5.858 & 6 & 0.439\\
4 & 2.057 & 1.245 &  5.506 & 5 & 0.357\\
\bottomrule
\end{tabular}
\caption{All-state X-only martingale audit for the surviving original
fixed-cash book plus spot.  Option means are basis points of initial spot.  No
horizon rejects the joint zero-mean null at 5\%.}
\label{tab:density-all-state}
\end{table}

The deterministic checks close much more tightly than Monte Carlo sampling:
the maximum density-tower residual is $6.49\times10^{-15}$, option-tower
residual $1.67\times10^{-16}$, spot-mean residual $2.60\times10^{-14}$,
four-roll decoder residual $1.77\times10^{-13}$, and first re-grid cash-price
residual $3.33\times10^{-16}$.

A matched run decomposed every weekly change into the fixed-expiry surface
branch and the spot-conditioning/re-grid component, using identical random
draws.  Direct inherited-cash and re-gridded valuations agree pathwise to at
most $2.22\times10^{-16}$ and have identical $z/J/p$ statistics.

\begin{table}[htbp]
\centering
\small
\begin{tabular}{rrrrr}
\toprule
Week & Fixed-expiry $p$ & Full first-member $p$ & Full pair $p$ & Max price difference\\
\midrule
1 & 0.366 & 0.684  & 0.139  & 0\\
2 & 0.538 & 0.163  & 0.128  & $2.22\times10^{-16}$\\
3 & 0.911 & 0.414  & 0.0208 & $2.22\times10^{-16}$\\
4 & 0.945 & 0.0808 & 0.546  & $2.22\times10^{-16}$\\
\bottomrule
\end{tabular}
\caption{Matched fixed-support and floating-coordinate decomposition.  The
first two $p$-value columns use one independent member of each pair; the third
uses antithetic pair means.}
\label{tab:density-regrid}
\end{table}

The week-3 pair-mean value is a nominal rejection at 5\%, but not at the
pre-specified four-look cutoff 0.0125.  It is also visibly sensitive to the
correlation eigenvalue cutoff: its $p$-value is 0.0429, 0.0208, 0.0879, and
0.295 at cutoffs 0.005, 0.01, 0.02, and 0.05.  Estimating covariance and the
mean on opposite half-samples gives 0.214 and 0.338.  No coordinate has
$|z|>2$ and the largest mean is 0.363bp.  Since the direct and re-gridded path
matrices are identical, this fluctuation cannot be caused by re-gridding.
The fixed-expiry branch alone has maximum option means of only 0.0046, 0.0056,
0.0013, and 0.0050bp over weeks 1--4.

Earlier development audits reached the same conclusion under a nonzero Black
backbone.  On 21 states with 10,000 paths per state, the four weekly maximum
coordinate $|z|$ values were 2.400, 3.203, 2.871, and 2.968, while the
covariance-aware Fisher-combined state-test $p$-values were 0.871, 0.672,
0.852, and 0.522.  A heterogeneous 10,000-path confirmation covering all 891
states gave unconditional maximum $|z|$ values 1.575, 1.970, 1.878, and 1.040
and $p$-values 0.702, 0.414, 0.067, and 0.317.  Its high-dimensional
state-feature-weighted values were 0.027, 0.759, 0.251, and 0.0089; they are
reported because they expose the poor calibration of hard spectral cutoffs in
small per-state samples, rather than being hidden behind the aggregate result.

\subsection{Repeated rolls, speed, and finite-support behavior}

A 16-week audit with the independent 10\% Black backbone used 10,000 paths and
all 891 states.  Across all rolls the largest decoder residual was
$2.35\times10^{-13}$, conditional mass and mean errors stayed below
$6.6\times10^{-14}$, and DLVsq remained strictly inside $(10^{-4},9)$.  The
maximum spot $|z|$ over 16 horizons was 2.126, the maximum rolling-option
$|z|$ was 2.417, and every rolling-option joint $p$-value was at least 0.0839.
With the last-expiry AR boundary, median cross-expiry dispersion of
$\sqrt{\mathrm{DLVsq}}$ was 0.0459 after week 1 and 0.0569 after week 16, so
the term structure did not collapse.  The corresponding implementation ran
160,000 path-week steps in 14.6 seconds.

The exact branch method is materially faster than a nonlinear finite-step
solve.  In the heterogeneous 10,000-path benchmark it processed 40,000
path-week steps in 8.84 seconds (4,527 per second), including 7.36 seconds to
construct 5,000 distinct initial kernels.  This was about 23 times faster than
the earlier finite-$h$ Newton transition (195 per second), but about 5.7 times
slower than the local SDE transition (25,649 per second).  The final X-only
audit, with more extensive all-state diagnostics, processed 400,000 path-week
steps in 169.8 seconds.  These timings are indicative rather than a controlled
hardware benchmark.

The remaining limitation is economic rather than a floating-point failure.  If
the lowest support node $x_0$ is revealed, then
\[
  \sum_rR^{b,j}_{0r}x_r=x_0,
  \qquad x_r\ge x_0,
\]
forces $R^{b,j}_{00}=1$.  The upper endpoint is analogous.  The $X$ law is then
absorbing, all surviving options are intrinsic, and interior DLVsq coordinates
are price-unidentified.  The positive-density formula $\eta=\psi/q$ no longer
applies at zero-mass nodes; any continued DLV coordinates there are an
economically invisible numerical convention.

In the four-week X-only run, boundary incidence was 0.003\%, 0.003\%, 0.003\%,
and 0.006\%.  In the 16-week run, the fully collapsed fraction grew from
0.01\% to 0.48\%; the median week-16 front density retained 23.44 effective
nodes, while its first percentile retained 1.15.  Avoiding this absorbing tail
requires dynamic support expansion, an explicit tail continuation, or a
boundary state carrying a conditional tail.  Re-gridding alone cannot create
support beyond an extreme node while preserving the inherited option surface.

The numerical evidence therefore supports three distinct statements: the
Heun implementation is a credible approximation to the continuous density
drift; the symmetric density construction is exactly martingale consistent on
the finite support; and weekly re-gridding preserves every surviving
fixed-cash price path by path.  It does not establish uniqueness, an
unbounded-support model, or the absence of boundary bias at horizons much
longer than those tested.
